\section*{Блок 5. Условная плотность и условное математическое ожидание}
\addcontentsline{toc}{section}{Блок 5. Условная плотность и условное математическое ожидание}

\subsection*{1. Теоретический фундамент и семантика преобразований}
\addcontentsline{toc}{subsection}{1. Теоретический фундамент и семантика преобразований}

Вычисление условных распределений абсолютно непрерывных случайных векторов при нелинейных связях является вычислительно ресурсоемкой задачей. Прямое применение определения условной вероятности через предел $ \lim_{\varepsilon \to 0} \P(X \in A \mid z - \varepsilon < Z < z + \varepsilon) $ на практике заменяется аналитическим аппаратом замены переменных.

Фундаментальная постановка: задан случайный вектор $(X, Y)$ с совместной плотностью $p_{X,Y}(x,y)$. В случае независимости $p_{X,Y}(x,y) = p_X(x)p_Y(y)$. Вводится новая случайная величина $Z = f(X,Y)$. Требуется определить $p_{X|Z}(x|z)$ и вычислить $\E[g(X,Y) \mid Z]$.

\subsection*{2. Универсальный аналитический аппарат: Метод Якобиана}
\addcontentsline{toc}{subsection}{2. Универсальный аналитический аппарат: Метод Якобиана}

Для построения совместной плотности $(X, Z)$ применяется диффеоморфизм областей значений. Вводится фиктивная переменная (чаще всего тождественная одной из исходных, например $U = X$). Строится обратное отображение: $X = U$, $Y = h(U, Z)$.

\begin{MethodBox}[Алгоритм вычисления условных характеристик]
Пусть требуется найти $\E[g(X,Y) \mid Z=z]$, где $Z = f(X,Y)$.

\begin{enumerate}[label=\textbf{\arabic*.}]
    \item \textbf{Построение совместной плотности.} Разрешить уравнение $z = f(x,y)$ относительно $y$: $y = h(x,z)$. Вычислить якобиан преобразования $J = \frac{\partial h(x,z)}{\partial z}$.
    Записать совместную плотность через исходную:
    \[
        p_{X,Z}(x,z) = p_{X,Y}(x, h(x,z)) \cdot |J|.
    \]
    
    \item \textbf{Строгий анализ носителя (Метод индикаторов).} Исходная плотность всегда содержит индикатор $I\{(x,y) \in D\}$. Выполнить подстановку $y = h(x,z)$ непосредственно в функцию-индикатор:
    \[
        I\{(x, y) \in D\} \implies I\{(x, h(x,z)) \in D\}.
    \]
    Разрешить полученную систему неравенств относительно $x$, получив границы интегрирования $A(z) < x < B(z)$, зависящие от параметра $z$.
    
    \item \textbf{Маргинализация.} Вычислить плотность условия $Z$, проинтегрировав по $x$:
    \[
        p_Z(z) = \int_{A(z)}^{B(z)} p_{X,Z}(x,z) \, dx.
    \]
    
    \item \textbf{Вычисление условной плотности и ожидания.} 
    \[
        p_{X|Z}(x|z) = \frac{p_{X,Z}(x,z)}{p_Z(z)} \cdot I\{A(z) < x < B(z)\}.
    \]
    \[
        \E[g(X,Y) \mid Z=z] = \int_{A(z)}^{B(z)} g(x, h(x,z)) p_{X|Z}(x|z) \, dx.
    \]
\end{enumerate}
\end{MethodBox}

\vspace{1em}

\subsection*{3. Продвинутые методы оптимизации вычислений}
\addcontentsline{toc}{subsection}{3. Продвинутые методы оптимизации вычислений}

Прямое интегрирование не всегда является оптимальным. Ряд задач решается в одну строку с использованием алгебраических свойств условного математического ожидания.

\begin{HackBox}[Свойство измеримости (Вынесение известного сомножителя)]
Относительно условия $Z$, любая функция $q(Z)$ ведет себя как детерминированная константа и выносится за знак математического ожидания:
\[
    \E[X \cdot q(Z) \mid Z] = q(Z) \cdot \E[X \mid Z].
\]
\textbf{Пример:} Требуется найти $\E[XY \mid Z]$, где $Z = Y/X \implies Y = XZ$.
Тогда $\E[XY \mid Z] = \E[X(XZ) \mid Z] = \E[X^2 Z \mid Z] = Z \cdot \E[X^2 \mid Z]$. 
Сложная функция двух переменных редуцируется к ожиданию квадрата одной переменной, что радикально упрощает интегрирование.
\end{HackBox}

\begin{HackBox}[Комбинаторная симметрия]
Если вектор $(X_1, \dots, X_n)$ состоит из независимых одинаково распределенных величин, а функция условия $S = \sum_{i=1}^n X_i$ симметрична относительно перестановок, то:
\[
    \E[X_1 \mid S] = \E[X_2 \mid S] = \dots = \E[X_n \mid S].
\]
В силу линейности условного ожидания: $\sum_{i=1}^n \E[X_i \mid S] = \E[S \mid S] = S$. Следовательно, $\E[X_k \mid S] = \frac{S}{n}$.
Данный метод позволяет вычислить ожидание мгновенно, минуя $n$-мерные интегралы.
\end{HackBox}

\begin{HackBox}[Обусловливание по немонотонным функциям]
При вычислении $\E[X \mid X^2]$, где $X \sim U(-a, a)$, прямое применение метода якобиана некорректно из-за отсутствия биекции.
При заданном $X^2 = z$, величина $X$ может принимать только два значения: $\sqrt{z}$ и $-\sqrt{z}$. В силу четности плотности исходного распределения, условное распределение дискретно и симметрично: $\P(X = \sqrt{z} \mid X^2=z) = \P(X = -\sqrt{z} \mid X^2=z) = \frac{1}{2}$. 
Отсюда $\E[X \mid X^2=z] = \sqrt{z} \cdot \frac{1}{2} - \sqrt{z} \cdot \frac{1}{2} = 0$. Строгий вывод: $\E[X \mid X^2] = 0$.
\end{HackBox}

\subsection*{4. Критические точки и топологические ошибки}
\addcontentsline{toc}{subsection}{4. Критические точки и топологические ошибки}

\begin{WarningBox}[Игнорирование динамических границ носителя]
Фундаментальная ошибка при поиске совместной плотности — использование фиксированных пределов интегрирования.
Пусть $X \in (0,1)$, $Y \in (0,1)$, $Z = Y/X$. Подстановка $Y = XZ$ в индикатор $I\{0 < Y < 1\}$ дает $0 < XZ < 1$.
Так как $Z > 0$ и $X > 0$, получаем $X < 1/Z$. Итоговая система для носителя переменной $X$:
\[
    0 < X < 1 \quad \text{и} \quad X < 1/Z \implies 0 < X < \min\left(1, \frac{1}{Z}\right).
\]
Таким образом, интеграл для маргинальной плотности $p_Z(z)$ распадается на две ветви: при $z \in (0,1]$ верхний предел равен 1, при $z \in (1, \infty)$ предел равен $1/z$. Интегрирование без учета функции минимума приводит к нарушению нормировки плотности.
\end{WarningBox}

\subsection*{5. Эвристики для экспресс-проверки результатов}
\addcontentsline{toc}{subsection}{5. Эвристики для экспресс-проверки результатов}

Аналитические расчеты плотностей подвержены арифметическим сбоям. Рекомендуется применять следующие инварианты для верификации итогового ответа $\E[g(X,Y) \mid Z=z] = \varphi(z)$:

\begin{enumerate}[label=\textbf{\arabic*.}]
    \item \textbf{Теорема о двойном ожидании (Формула полного математического ожидания).} 
    Безусловное математическое ожидание должно восстанавливаться усреднением условного:
    \[
        \E[\varphi(Z)] = \int_{-\infty}^{\infty} \varphi(z) p_Z(z) \, dz \equiv \E[g(X,Y)].
    \]
    Обе части вычисляются независимо. Совпадение является строгим критерием корректности найденной условной плотности и самого ожидания.
    
    \item \textbf{Принцип сохранения границ (Ограниченность).}
    Если исходная функция ограничена почти наверное, например $m \le g(X,Y) \le M$, то и полученная функция условного ожидания $\varphi(z)$ обязана лежать в пределах $[m, M]$ для всех допустимых значений $z$. Выход за границы указывает на неверно определенный носитель или потерянный якобиан.
    
    \item \textbf{Асимптотика экстремальных условий.}
    Рассмотрение предельных значений $z$ позволяет выявить топологические ошибки. Например, для отношения $Z = Y/X$ при положительных $X$ и $Y$: предел $z \to 0$ означает, что $Y \ll X$. Условная плотность $X$ в этом пределе должна концентрироваться около своих максимумов. Если $\lim_{z \to 0} \E[X \mid Z=z]$ не согласуется с физическим смыслом величин, в расчете допущена алгебраическая ошибка.
\end{enumerate}
